%!TEX root = ../Master_Thesis.tex

\chapter{Conclusion}

This chapter will wrap up the discussion on the interpreter and it's various implementation approaches.
A summary will be given, that will outline the key aspects of all of the previous chapter.
After that, the evaluation result will be discussed and the research questions of section \ref{sec:goals} will be answered.
The current state of art will then be described using other examples.
This thesis will then end with a list of next steps resulting from the work and references to related papers and work.

\section{Summary}


Fabrik19 is a company, that develops marketing and e-commerce applications for digital signage and mobile apps.
To be able to update the contents, layouts and behaviors of the applications, a model interpreter infrastructure has been developed, so there is no need to update the local application bundles on the physical signage devices or the mobile applications.

The model defines layout, behavior and data.
The configuration allows the definition of complex layouts and interaction.
This makes it very powerful and functional, while still being very flexible.
If the configuration changes, the applications download the changes and updates it's behavior accordingly.

Because of low starting requirements, the interpreter's performance is noticeable slower than a native application's performance, that implements the functionality directly and natively.
This is because a software architecture has been implemented, that directly maps the models capabilities to native \ac{UI} elements.
Since changes to the \ac{UI} are happening multiple times with only one user interaction, the \ac{UI} thread is very busy.
An interaction example is a database change, that could trigger layout updates.
If the interaction produces multiple interactions, all the changes are applied to the \ac{UI}, blocking the main thread.

An alternative approach of implementing the interpreter is the virtual machine approach, which simulates the app's runtime with basic data structures.
This system is independent to any native \ac{UI} implementation and the main thread.
This means, that any change can be evaluated and the virtual \ac{UI} hierarchy can be created on custom background threads.
After the evaluation is finished and every calculation is done, the changes are applied to the actual native \ac{UI}.
The application of changes is fast, because only the resulting changes are being applied.
This means there is no duplicated work.
The idea is inspired by the virtual DOM of React Native and the event zoning of Angular.




\section{Discussion}



\section{State of the Art}

\section{Next Steps}

This section will discuss further development approaches and the development steps that should now be taken by the team.
This includes alternative architectures and concrete implementation approaches.

\subsection{Continuous Implementation}

The next step is the actual implementation of the interpreter.
It could be considered a full rewrite, but the more intelligent solution is a continuous evolution, which transfers the application into the target architecture.
The approach that was described fully separates the virtual runtime and the user interface.
This radical step can be applied to the interpreter by separating all the aspects step by step.

In the beginning, the UI behavior can be extracted from the various layer implementations, by sub-classing native view elements and implement view specific behavior.
This way, the dirt flushing mechanism can be implemented step by step.
The next major problems are the view controller, which are the windows that present the node layouts.
Those also need to be separated, so the window management is working with the virtual windows, which directly use the native elements.
After that, the dirt flush mechanism is fully implemented so that every change only results in a dirt flag, which will then be applied to the view after the interactions are finished.

Each step can be considered a sprint between which the application is executable.
Every change is a step in the right direction, because a good abstraction supports maintainability and extensibility.
After every changes, the team can evaluated if the next step is even possible or if the architecture is not suitable at all.

\subsection{Hybrid Approaches}

\subsection{Runtime Models}

\subsection{Abstraction}


\section{Related Work}